\section{价值链热力图}

产业链分析的核心不是把所有相关公司都列出来，而是判断价值池向哪里迁移。AI服务器把PCB从传统制造件推向高速互联件，价值从普通多层板迁移到高层数PCB、HDI、高频高速CCL、载板材料和精密设备耗材。这个迁移不均匀：越靠近高速信号完整性、客户认证和良率瓶颈，越可能获得定价权；越靠近通用产能，越容易在扩产周期中被价格竞争稀释。

因此，本章先把产业链拆成可投资环节，再讨论平台链条和公司关系。表格只是为了压缩比较维度，真正的结论是：高频高速CCL和高端AI PCB是第一顺位；载板/封装基板和设备耗材是第二顺位；上游通用材料和低端PCB不应被简单纳入AI主线。

\begin{exhibitbox}[图表8：按瓶颈和可投资性划分的环节热力图]
\small
\begin{tabularx}{\textwidth}{L{2.4cm}C{1.8cm}C{1.8cm}C{1.8cm}C{1.8cm}X}
\toprule
\textbf{环节} & \textbf{AI需求} & \textbf{定价} & \textbf{认证壁垒} & \textbf{证据} & \textbf{代表公司} \\
\midrule
上游材料 & \cellcolor{yellow!18}中 & \cellcolor{yellow!18}中 & \cellcolor{yellow!18}中 & \cellcolor{yellow!18}E3 & 中国巨石、中材科技、国际复材、宏和科技 \\
高频高速CCL & \cellcolor{green!18}高 & \cellcolor{green!18}高 & \cellcolor{green!18}高 & \cellcolor{green!18}E2 & 生益科技、华正新材、南亚新材 \\
PCB制造 & \cellcolor{green!18}高 & \cellcolor{yellow!18}中 & \cellcolor{green!18}高 & \cellcolor{green!18}E2 & 沪电股份、胜宏科技、深南电路、鹏鼎控股 \\
IC载板 & \cellcolor{yellow!18}中 & \cellcolor{yellow!18}中 & \cellcolor{green!18}高 & \cellcolor{yellow!18}E2/E3 & 深南电路、兴森科技、华正新材 \\
设备 & \cellcolor{yellow!18}中 & \cellcolor{yellow!18}中 & \cellcolor{yellow!18}中 & \cellcolor{yellow!18}E3 & 大族数控、芯碁微装、东威科技$^\dagger$、凯格精机 \\
耗材 & \cellcolor{yellow!18}中 & \cellcolor{yellow!18}中 & \cellcolor{yellow!18}中 & \cellcolor{yellow!18}E3 & 鼎泰高科、中钨高新 \\
\bottomrule
\end{tabularx}
\sourcenote{证据等级对应 ch02 图表2-1 的 E1--E4 体系：E2=原始券商报告/明确归因转载；E2/E3=部分券商加部分跨来源推断；E3=跨来源产业逻辑。高频高速CCL行：CMBI行业结构证据 + 材料/价格序列；PCB制造行：H股官方申请资料环节摘要；IC载板行：券商封装基板报告 + 跨来源推断；上游材料/设备/耗材行：跨来源产业逻辑。$\dagger$ 东威科技非本报告覆盖标的。}
\end{exhibitbox}

从热力图看，需求强并不自动等于投资价值高。高端PCB和CCL的需求确定性最高，但估值也最拥挤；设备耗材有较强二阶弹性，但订单节奏滞后且波动更大；IC载板需要看良率、折旧和客户认证，不能只按AI算力主题给估值。

\begin{exhibitbox}[图表8b：CMBI/Prismark行业结构交叉验证]
\footnotesize
\begin{tabularx}{\textwidth}{L{2.6cm}X X}
\toprule
\textbf{维度} & \textbf{机构证据} & \textbf{报告含义} \\
\midrule
全球PCB周期 & CMBI引用Prismark数据，2025E全球PCB增速由7.6\%上修至12.8\%；2024年销售额约740亿美元。 & AI基础设施同时是周期和结构需求驱动，但增长假设仍需订单验证。 \\
高端结构 & 18层以上多层板、HDI和载板预计为突出环节；中国约占2024年全球PCB产值56\%。 & 支持聚焦高端A股供应商，而不是泛低端PCB暴露。 \\
CCL集中度 & 全球CCL前十大供应商2024年占比约77\%；高速专用CCL同比增长约50\%。 & 解释材料通胀时CCL龙头可能保留定价能力。 \\
成本传导 & 铜被引用为PCB原材料成本的60\%--70\%，高端AI HDI毛利显著高于标准PCB。 & 强化本报告对高端结构赢家和大宗成本风险的区分。 \\
\bottomrule
\end{tabularx}
\sourcenote{CMBI行业结构证据；券商二季度行业更新口径，不披露具名客户收入}
\end{exhibitbox}

行业结构数据进一步支持“高端优先”的判断。全球PCB增长被AI基础设施上修，但普通PCB和高端AI PCB的盈利含义完全不同：前者更多体现周期修复，后者体现材料等级、客户认证和良率壁垒。报告后续给沪电、胜宏、深南、生益更高关注度，正是因为它们更接近这些瓶颈。

\section{平台链条图}

平台链条决定需求从哪里来。NVIDIA/Rubin链条强调高层高速PCB和M9材料，ASIC/TPU链条强调HDI和高速服务器板，国产算力链条更强调国产高端材料和载板认证，光模块/网络链条则更看重交换机板、光模块PCB和低损耗材料。不同平台对应不同公司，不应把“AI服务器”作为一个笼统变量。

\begin{exhibitbox}[图表9：需求平台到A股暴露]
\small
\begin{tabularx}{\textwidth}{L{3cm}X X L{2.5cm}}
\toprule
\textbf{需求节点} & \textbf{所需PCB/CCL能力} & \textbf{可能受益A股} & \textbf{置信度} \\
\midrule
NVIDIA GB300/Rubin & 高层高速PCB、背板、中板、M9材料、CoWoP相关结构 & 沪电股份、胜宏科技、深南电路、生益科技 & \evidencetag{券商陈述} \\
北美ASIC & HDI、高速CCL、服务器板、交换机板 & 胜宏科技、沪电股份、生益科技 & \evidencetag{券商陈述} \\
国产算力 & 国产高端CCL、CBF/BT材料、PCB与载板认证 & 华正新材、深南电路、生益科技、兴森科技 & \evidencetag{券商陈述} \\
光模块/网络 & 光模块PCB、交换机PCB、高速材料 & 沪电股份、深南电路、景旺电子等 & \evidencetag{推断} \\
\bottomrule
\end{tabularx}
\sourcenote{券商陈述：平台-公司映射来自归档券商报告（高盛沪电AI PCB、JPM胜宏AI PCB、CMBI行业结构证据）；非官方客户收入拆分，具名平台份额未经官方披露}
\end{exhibitbox}

这里最需要警惕的是证据强弱差异。平台链条可以解释为什么某些公司进入观察范围，但除非有官方披露或原始报告支持，不能把具体客户收入写成确定事实。后续估值表对胜宏、鹏鼎和华正保留更宽区间，就是因为平台链条的收入拆分仍不完整。

\section{关系矩阵}

公司关系矩阵把产业链从“环节”推进到“标的”。同处PCB链条的公司，投资含义不同：PCB制造看订单和良率，CCL材料看价格和代际认证，设备耗材看扩产节奏，载板材料看研发和量产转换。

\begin{dashboardbox}[公司关系置信度]
\small
\begin{tabularx}{\textwidth}{L{2.3cm}L{3.0cm}X L{2.5cm}}
\toprule
\textbf{公司} & \textbf{角色} & \textbf{投资含义} & \textbf{标签} \\
\midrule
沪电股份 & AI PCB/背板 & 直接受益规格升级，估值受订单节奏影响。 & \evidencetag{券商陈述} \\
胜宏科技 & AI PCB/HDI & 弹性高，客户集中度和A/H可比性需验证。 & \evidencetag{券商陈述} \\
深南电路 & PCB+IC载板 & 财务交付证据较强，双轮驱动。 & \evidencetag{券商陈述} \\
生益科技 & 高频高速CCL & 材料涨价和结构升级主线。 & \evidencetag{券商陈述} \\
华正新材 & 高速CCL/CBF/BT & 国产算力和材料期权，估值对兑现敏感。 & \evidencetag{券商陈述} \\
设备耗材组 & 钻孔/LDI/电镀/钻针 & 扩产二阶弹性，订单波动高。 & \evidencetag{推断} \\
\bottomrule
\end{tabularx}
\end{dashboardbox}

\section{供应商经济性}

供应商经济性解释毛利率能否持续。AI产品放量通常伴随更高等级铜箔、覆铜板、半固化片和设备投入，短期可能推高收入和ASP，但也会提高材料采购集中度和营运资本占用。只有当客户认证、价格传导和供应链稳定同时成立，高端产品结构才会真正转化为利润率。

\begin{exhibitbox}[图表9b：H股文件供应商集中度与材料成本占比]
\footnotesize
\begin{tabularx}{\textwidth}{L{1.8cm}X X}
\toprule
\textbf{公司} & \textbf{供应商证据} & \textbf{模型含义} \\
\midrule
沪电股份 & 原材料成本占销售成本比例从2023年55.8\%升至2026Q1的62.3\%；五大供应商采购占比约41\%--42\%，最大供应商约21\%。 & 铜箔、覆铜板、半固化片成本传导是毛利率关键。 \\
沪电股份 & 主要供应商多为覆铜板及半固化片供应商，信用期60--150天；设备供应商信用期约30天。 & 供应商信用期支持DPO和现金转换假设。 \\
胜宏科技 & 原材料成本占销售成本比例从2023年58.6\%升至2025年65.9\%；五大供应商采购占比2023年31.4\%、2024年降至29.3\%、2025年跃升至45.2\%。 & 2025年集中度大幅跃升，与高端HDI/MLPCB放量伴随的材料和设备集中度提升一致。 \\
胜宏科技 & 供应商A为韩国电子材料供应商，供应商B为广东高端CCL开发商，供应商H为激光智能装备公司。 & 上游材料和设备能力影响产能爬坡与毛利率。 \\
\bottomrule
\end{tabularx}
\sourcenote{官方文件：H股申请资料供应商与原材料披露}
\end{exhibitbox}

供应商数据对投资判断的含义是：不要只看收入增速。胜宏科技这类高弹性公司如果材料成本和设备投入同步上升，利润率能否保持才是关键；沪电股份的供应商信用期更长、供应商集中度结构更稳定，毛利率可持续性证据更强。该判断在 ch08 估值中体现为沪电采用更窄的合理价值区间，胜宏、鹏鼎与华正保留更宽区间（见本章平台链条图后的边界说明）。
